\chapter{Hipoteza continuum i zastosowania wymuszania}
\label{ch:forcing-continuum}
\chaptermark{Wymuszanie i hipoteza continuum}

\section{Cel i zakres argumentu}
Hipoteza continuum (CH) jest zdaniem
$2^{\aleph_0}=\aleph_1$. Równoważnie:
nie istnieje moc ściśle pomiędzy
$\aleph_0$ a mocą $\mathcal P(\omega)$.
W poprzednim rozdziale uzasadniliśmy,
że dla przeliczalnego tranzytywnego
modelu zbiorowego $M$ teorii ZFC,
pojęcia wymuszania $\mathbb P\in M$
i filtru $M$-generycznego $G$,
rozszerzenie $M[G]$ również spełnia ZFC.

Pokażemy dwie odmienne konstrukcje.
Pierwsza dodaje wiele ciągów binarnych
i zachowuje liczby kardynalne modelu
wyjściowego, uzyskując $\neg\mathrm{CH}$.
Druga zmienia moce większych zbiorów
bez dodawania nowych podzbiorów $\omega$,
co pozwala uzyskać CH.

\begin{uwaga}[Perspektywa wewnętrzna i zewnętrzna]
\label{rem:forcing34-perspective}
W całym rozdziale $M$ jest
\emph{przeliczalnym tranzytywnym modelem
zbiorowym ZFC}. Własności
takie jak „nieprzeliczalny”,
„$\omega_1$”, „$\omega_2$” i
„warunek przeliczalny”,
jeśli dotyczą obiektów $M$,
są rozumiane \emph{wewnątrz $M$}.
Przykładowo $\omega_2^M$ jest
liczbą porządkową należącą do $M$,
choć zewnętrznie jest przeliczalna.
Warunki przeliczalności modelu $M$
używamy wyłącznie zewnętrznie, aby
otrzymać filtr generyczny.
Istnienia takiego $M$ nie wyprowadzamy
z samej $\operatorname{Con}(\mathrm{ZFC})$.
\end{uwaga}

\section{Liczby porządkowe w rozszerzeniu}

\begin{lemat}[Wymuszanie nie dodaje liczb porządkowych]
\label{lem:forcing34-no-ordinals}
Niech $(\mathbb P,\leq,1_{\mathbb P})\in M$
będzie pojęciem wymuszania, a
$G\subseteq\mathbb P$ filtrem
$M$-generycznym. Wtedy
\[
 \operatorname{Ord}^{M[G]}
   =\operatorname{Ord}^M,
\]
przy czym równość oznacza, że
każda liczba porządkowa należąca
do $M[G]$ należy do $M$ i odwrotnie.
\end{lemat}
\begin{proof}
Włączenie od prawej do lewej jest
natychmiastowe, bo $M\subseteq M[G]$
i oba modele są tranzytywne.
Dla drugiego włączenia wybierzmy
nazwę $\tau\in M$ o wartości
$\tau^G=\alpha$, gdzie $\alpha$
jest liczbą porządkową w $M[G]$.
Z hierarchii nazw poprzedniego
rozdziału otrzymujemy w $M$
liczbę porządkową $\gamma$
ograniczającą rangi nazw
występujących dziedzicznie w $\tau$.
Indukcja po tej hierarchii pokazuje
w otaczającym uniwersum, że
$\operatorname{rank}(\tau^G)\leq\gamma$:
każdy element $\tau^G$ jest wartością
nazwy o mniejszej randze.
Ponieważ $\alpha$ jest liczbą
porządkową, $\operatorname{rank}(\alpha)
=\alpha$, a więc $\alpha\leq\gamma$.
Tranzytywność $M$ i fakt
$\gamma\in M$ dają $\alpha\in M$
(również gdy $\alpha=\gamma$).
Nie otrzymujemy nowych liczb
porządkowych, choć ich \emph{moce}
mogą się zmieniać.
\end{proof}

\section{Warunek przeliczalnych antyłańcuchów}

\begin{definicja}[Antyłańcuch i warunek ccc]
\label{def:forcing34-ccc}
Niech $(\mathbb P,\leq)$ będzie
niepustym pojęciem wymuszania.
Zbiór $A\subseteq\mathbb P$ nazywamy
\emph{antyłańcuchem}, gdy każde dwa
różne elementy $A$ są niezgodne,
czyli nie mają wspólnego wzmocnienia.
Powiemy, że $\mathbb P$ spełnia
\emph{warunek przeliczalnych
antyłańcuchów} (ccc), gdy każdy
jego antyłańcuch jest co najwyżej
przeliczalny. W zastosowaniu do
$\mathbb P\in M$ wymagamy
$M\models\text{„$\mathbb P$ spełnia ccc”}$.
\end{definicja}

\begin{definicja}[System $\Delta$]
\label{def:forcing34-delta-system}
Niech $I$ będzie zbiorem, a
$\mathcal F$ rodziną jego skończonych
podzbiorów. Rodzina $\mathcal F$
jest \emph{systemem $\Delta$ o rdzeniu $R$},
gdy $R\subseteq I$ oraz dla
wszystkich różnych $A,B\in\mathcal F$
mamy $A\cap B=R$.
\end{definicja}

\begin{lemat}[Lemat o systemie $\Delta$]
\label{lem:forcing34-delta}
W ZFC każda nieprzeliczalna rodzina
różnych skończonych zbiorów
zawiera nieprzeliczalną podrodzinę,
która jest systemem $\Delta$.
\end{lemat}
\begin{proof}
Z dowolnej nieprzeliczalnej rodziny
wybieramy podrodzinę mocy $\aleph_1$.
Ponieważ wszystkie jej elementy
są skończone, z regularności
$\aleph_1$ istnieje $n\in\omega$
oraz podrodzina mocy $\aleph_1$
złożona z $n$-elementowych zbiorów.
Indukujemy po $n$.

Dla $n=0$ twierdzenie jest oczywiste;
rodzina różnych zbiorów pustych nie
ma zresztą mocy $\aleph_1$.
Załóżmy, że $n>0$. Jeżeli pewien
element $x$ należy do $\aleph_1$
wielu zbiorów rodziny, usuwamy
$x$ z każdego z nich. Dostajemy
$\aleph_1$ różnych zbiorów o
$n-1$ elementach. Hipoteza
indukcyjna daje system $\Delta$
o rdzeniu $R$, a po ponownym
dołączeniu $x$ otrzymujemy
system o rdzeniu $R\cup\{x\}$.

W przeciwnym przypadku każdy element
należy do co najwyżej przeliczalnie
wielu zbiorów rozpatrywanej rodziny.
Rekurencyjnie przez wszystkie
$\xi<\omega_1$ wybierzemy parami
rozłączne zbiory $F_\xi$.
Przed krokiem $\xi$ wybrano
przeliczalnie wiele zbiorów,
gdyż $\xi<\omega_1$ jest
liczbą porządkową przeliczalną.
Ich suma jest przeliczalna;
każdy jej element występuje
w co najwyżej przeliczalnie
wielu zbiorach rodziny.
Zatem tylko przeliczalnie wiele
zbiorów jest wykluczonych i
można wybrać następny rozłączny.
Otrzymana rodzina mocy $\aleph_1$
jest systemem $\Delta$ o pustym
rdzeniu. To kończy indukcję.
\end{proof}

\begin{definicja}[Dodawanie wielu ciągów Cohena]
\label{def:forcing34-many-cohen}
Niech $\kappa\in M$ będzie niezerową
liczbą kardynalną modelu $M$.
Definiujemy w $M$ pojęcie wymuszania
\[
 \operatorname{Add}(\omega,\kappa)^M
 =\{p:p\text{ jest skończoną funkcją }
       \kappa\times\omega\longrightarrow2\}.
\]
Porządek jest odwrotną inkluzją:
$q\leq p$ wtedy i tylko wtedy,
gdy $q\supseteq p$. Największym
warunkiem jest funkcja pusta.
\end{definicja}

\begin{lemat}[Warunek ccc dla wielu ciągów Cohena]
\label{lem:forcing34-add-ccc}
Niech $M$ będzie tranzytywnym modelem
ZFC, a $\kappa$ niezerową liczbą
kardynalną należącą do $M$.
Wówczas
\[
 M\models
 \text{„$\operatorname{Add}(\omega,\kappa)$
       spełnia ccc”}.
\]
\end{lemat}
\begin{proof}
Prowadzimy cały argument wewnątrz $M$.
Gdyby istniał nieprzeliczalny
antyłańcuch $A$, wybralibyśmy
$\aleph_1$ różnych warunków $p$
należących do $A$.
Ich dziedziny są skończonymi
podzbiorami $\kappa\times\omega$.
Ponieważ dla jednej ustalonej
dziedziny istnieje tylko skończenie
wiele funkcji do $\{0,1\}$,
zbiór wybranych dziedzin jest
nieprzeliczalny. Lemat
\ref{lem:forcing34-delta} daje
nieprzeliczalną podrodzinę
o parami jednakowych przecięciach
dziedzin równych skończonemu
rdzeniowi $R$.
Na $R$ istnieje tylko $2^{|R|}$
możliwych przypisań wartości.
Wybieramy dwie różne funkcje
$p,q$ z tej podrodziny,
które zgadzają się na $R$.
Ponieważ ich dziedziny poza $R$
są rozłączne, $p\cup q$ jest
funkcją i wspólnym wzmocnieniem
$p$ i $q$. Przeczy to temu,
że $A$ jest antyłańcuchem.
\end{proof}

\begin{przyklad}[Zgodność pomimo różnych dziedzin]
\label{ex:forcing34-compatible}
Niech $\kappa$ ma co najmniej
trzy elementy i niech
\[
\begin{aligned}
p&=\{((0,0),1),((1,2),0)\},\\
q&=\{((1,2),0),((2,1),1)\}.
\end{aligned}
\]
Wspólny rdzeń dziedzin to
$\{(1,2)\}$, gdzie oba warunki
przypisują $0$. Dlatego
$p\cup q$ jest ich wspólnym
wzmocnieniem. Gdyby natomiast
$q(1,2)=1$, warunki byłyby
niezgodne. Nie chodzi więc o
samą wielkość dziedzin, lecz
o zgodność funkcji na przecięciu.
\end{przyklad}

\section{Dlaczego ccc zachowuje moce}

\begin{lemat}[Gęste rozstrzyganie wartości porządkowej]
\label{lem:forcing34-decide-ordinal}
Niech $\mathbb P\in M$ będzie pojęciem
wymuszania, $\dot f\in M$ nazwą,
$\beta,\lambda\in M$ liczbami
porządkowymi oraz $p\in\mathbb P$
warunkiem wymuszającym, że
$\dot f$ jest funkcją $\check\beta$
w $\check\lambda$. Dla każdego
$\xi<\beta$ zbiór
\[
 D_\xi=\{q\leq p:\exists\eta<\lambda\
      q\Vdash\dot f(\check\xi)=\check\eta\}
\]
jest gęsty poniżej $p$ i należy
do $M$.
\end{lemat}
\begin{proof}
Jest to podzbiór $\mathbb P$
definiowalny w $M$ przy użyciu
Twierdzenia~\ref{thm:forcing33-definability},
a więc należy do $M$.
Gdyby nie był gęsty poniżej $p$,
istniałby $q\leq p$, poniżej
którego żaden warunek nie wymusza
równości z żadną $\check\eta$.
Z konstrukcji filtru generycznego
w Rozdziale~\ref{ch:forcing-foundations}
wybierzmy filtr $H$ zawierający $q$.
Skoro $p\in H$, w $M[H]$
wartość $\dot f^H(\xi)$ jest
elementem $\lambda$, a zatem
równym pewnej $\eta<\lambda$
należącej do $M$. Lemat prawdy
daje $r\in H$ wymuszające
$\dot f(\check\xi)=\check\eta$.
Wspólne wzmocnienie $r$ i $q$
należałoby do $D_\xi$,
sprzecznie z wyborem $q$.
\end{proof}

\begin{twierdzenie}[Zachowanie liczb kardynalnych przez ccc]
\label{thm:forcing34-ccc-preserves}
Niech $M$ będzie przeliczalnym
tranzytywnym modelem ZFC,
a $\mathbb P\in M$ pojęciem
wymuszania spełniającym ccc
w sensie $M$. Dla każdego filtru
$M$-generycznego $G$ każda
liczba kardynalna modelu $M$
pozostaje liczbą kardynalną
modelu $M[G]$. W szczególności
$\omega_1^{M[G]}=\omega_1^M$
i $\omega_2^{M[G]}=\omega_2^M$.
\end{twierdzenie}
\begin{proof}
Z Lematu~\ref{lem:forcing34-no-ordinals}
wiemy, że $M[G]$ i $M$
mają te same liczby porządkowe.
Załóżmy, że pewna nieprzeliczalna
liczba kardynalna $\lambda$ z $M$
przestała być kardynalna w $M[G]$.
Wtedy dla pewnego $\beta<\lambda$
istnieje w $M[G]$ surjekcja
$f:\beta\to\lambda$.
Wybierzmy nazwę $\dot f\in M$
tej funkcji i warunek $p\in G$
wymuszający, że $\dot f$ jest
taką surjekcją.

Dla każdego $\xi<\beta$
rozpatrujemy w $M$ zbiór
$D_\xi$ z Lematu~\ref{lem:forcing34-decide-ordinal}.
Jest gęsty poniżej $p$,
więc na mocy aksjomatu wyboru
można wybrać maksymalny
antyłańcuch $A_\xi\subseteq D_\xi$,
\emph{maksymalny poniżej $p$}.
Jego istnienie uzasadnia się
lematem Kuratowskiego--Zorna:
porządek na antyłańcuchach
przez inkluzję ma górne
ograniczenia dla łańcuchów.
Wobec ccc każdy $A_\xi$
jest co najwyżej przeliczalny
\emph{w $M$}. Niech
\[
 S_\xi=\{\eta<\lambda:
  \exists q\in A_\xi\
  q\Vdash\dot f(\check\xi)=\check\eta\}.
\]
Zatem $S_\xi\in M$ jest co najwyżej
przeliczalny. Zastępowanie w $M$
pozwala utworzyć wszystkie
zbiory $S_\xi$ jednocześnie.
W konsekwencji
$S=\bigcup_{\xi<\beta}S_\xi$
spełnia w $M$ oszacowanie
\[
 |S|\leq\max(\aleph_0,|\beta|)<\lambda.
\]

Dlaczego w każdym rozszerzeniu
generycznym zawierającym $p$
mamy $f(\xi)\in S_\xi$?
Maksymalność $A_\xi$ poniżej
$p$ zapewnia, że zbiór
\[
 E_\xi=\{q:\exists a\in A_\xi\
              q\leq a\}
       \cup\{q:q\perp p\}
\]
jest gęsty w całym $\mathbb P$
i należy do $M$.
Jeżeli $p\in H$ dla filtru
$M$-generycznego $H$,
to $H$ przecina $E_\xi$,
a ponieważ żaden warunek
$H$ nie jest niezgodny z $p$,
filtr $H$ zawiera pewien
$a\in A_\xi$. Warunek $a$
ustala wartość $f^H(\xi)$
w zbiorze $S_\xi$.
Zatem
$\operatorname{ran}f^H\subseteq S$,
wbrew temu, że $f^H$ ma być
surjekcją na $\lambda$.
Sprzeczność kończy dowód.

Liczba $\omega$ pozostaje
kardynalna także w każdym
tranzytywnym modelu ZFC
o tych samych liczbach porządkowych.
Ponieważ wszystkie kardynały
modelu $M$ są zachowane,
jego kolejne kardynały
$\omega_1$ i $\omega_2$
mają tę samą wartość
w $M[G]$.
\end{proof}

\section{Wymuszanie zaprzeczenia CH}

\begin{twierdzenie}[Wiele różnych nowych ciągów]
\label{thm:forcing34-many-reals}
Niech $\kappa\in M$ będzie
nieskończoną liczbą kardynalną,
$\mathbb P=\operatorname{Add}(\omega,\kappa)^M$,
a $G$ filtrem $M$-generycznym.
Wówczas funkcja
$F=\bigcup G$ ma dziedzinę
$\kappa\times\omega$ i wartości
w $\{0,1\}$. Dla każdego
$\alpha<\kappa$ funkcja
$c_\alpha(n)=F(\alpha,n)$
jest elementem $2^\omega\cap M[G]$.
Ponadto
$c_\alpha\neq c_\beta$
dla różnych $\alpha,\beta<\kappa$.
\end{twierdzenie}
\begin{proof}
Skierowanie $G$ zapewnia, że
każde dwa warunki należące
do $G$ są zgodnymi funkcjami.
Dlatego ich suma $F=\bigcup G$
jest funkcją i należy do $M[G]$.
Dla każdej pary $(\alpha,n)$
zbiór warunków określonych
na tej współrzędnej jest
gęsty w $\mathbb P$ i należy
do $M$. Filtr $G$ go przecina,
więc $F$ jest całkowitą funkcją
na $\kappa\times\omega$.

Ustalmy różne $\alpha,\beta<\kappa$.
W modelu $M$ rozpatrzmy
\[
\begin{aligned}
 D_{\alpha,\beta}
  =\{p\in\mathbb P:\exists n<\omega\
  (&(\alpha,n),(\beta,n)
            \in\operatorname{dom}p\\
   &{}\land p(\alpha,n)\ne p(\beta,n))\}.
\end{aligned}
\]
Zbiór jest gęsty: dla dowolnego
skończonego $p$ wybieramy
$n$, dla którego obie współrzędne
nie należą do $\operatorname{dom}p$,
i przedłużamy $p$ o wartości
$0$ i $1$ na tych współrzędnych.
Ponieważ $D_{\alpha,\beta}\in M$,
mamy $G\cap D_{\alpha,\beta}\ne\varnothing$.
Wtedy dla wskazanego $n$
wartości $c_\alpha(n)$ i
$c_\beta(n)$ są różne.
Każde $c_\alpha$ jest podzbiorem
$\omega$ kodowanym przez funkcję
binarną, zatem należy do $M[G]$.
\end{proof}

\begin{twierdzenie}[Generyczny model zaprzeczający CH]
\label{thm:forcing34-not-ch}
Niech $M$ będzie przeliczalnym
tranzytywnym modelem zbiorowym ZFC.
Połóżmy
$\kappa=\omega_2^M$ i
$\mathbb P=\operatorname{Add}(\omega,\kappa)^M$.
Dla każdego filtru $M$-generycznego
$G\subseteq\mathbb P$ zachodzi
\[
 M[G]\models\mathrm{ZFC}+\neg\mathrm{CH}.
\]
Nie trzeba zakładać $M\models\mathrm{CH}$.
\end{twierdzenie}
\begin{proof}
Twierdzenie~\ref{thm:forcing33-zfc}
daje $M[G]\models\mathrm{ZFC}$.
Z Lematu~\ref{lem:forcing34-add-ccc}
i Twierdzenia~\ref{thm:forcing34-ccc-preserves}
wynika, że $\kappa$ pozostaje
$\omega_2$ w modelu $M[G]$.
Natomiast Twierdzenie
\ref{thm:forcing34-many-reals}
daje w $M[G]$ iniekcję
$\alpha\mapsto c_\alpha$
ze zbioru $\kappa$ do $2^\omega$.
Stąd wewnątrz $M[G]$ mamy
\[
 2^{\aleph_0}\geq\aleph_2>\aleph_1.
\]
To jest zaprzeczenie CH.
\end{proof}

\section{Przeliczalna domkniętość i brak nowych ciągów}

\begin{definicja}[Przeliczalna domkniętość]
\label{def:forcing34-countably-closed}
Pojęcie wymuszania $(\mathbb P,\leq)$
jest \emph{przeliczalnie domknięte},
jeżeli dla każdego ciągu
$p_0\geq p_1\geq\cdots$
warunków z $\mathbb P$
istnieje $q\in\mathbb P$
takie, że $q\leq p_n$
dla każdego $n\in\omega$.
Jeśli $\mathbb P\in M$, to
wymagamy tej własności dla
ciągów należących do $M$,
czyli w sensie
$M\models\text{„$\mathbb P$
jest przeliczalnie domknięte”}$.
\end{definicja}

\begin{lemat}[Domkniętość nie dodaje ciągów liczb porządkowych]
\label{lem:forcing34-closed-no-sequences}
Niech $\mathbb P\in M$ będzie
przeliczalnie domkniętym pojęciem
wymuszania, a $G$ filtrem
$M$-generycznym. Jeśli
$f\in M[G]$ jest funkcją
$f:\omega\to\operatorname{Ord}^{M[G]}$,
to $f\in M$. W szczególności
$M[G]$ i $M$ mają dokładnie
te same podzbiory $\omega$
i tę samą liczbę porządkową
$\omega_1$.
\end{lemat}
\begin{proof}
Wybierzmy nazwę $\dot f\in M$
oraz $p\in G$ wymuszające, że
$\dot f$ jest funkcją z $\omega$
w klasę liczb porządkowych.
Z Lematu~\ref{lem:forcing34-no-ordinals}
wiemy, że jej wartości są
liczbami porządkowymi należącymi do $M$.

Dla każdego $n\in\omega$ zbiór
warunków poniżej $p$ ustalających
wartość $\dot f(\check n)$
jako pewną liczbę porządkową
z $M$ jest gęsty poniżej $p$.
Dowód jest identyczny z argumentem
w Lemacie~\ref{lem:forcing34-decide-ordinal},
z tą różnicą, że kwantyfikator
po możliwych liczbach porządkowych
nie ma ustalonego ograniczenia,
lecz jest zwykłym kwantyfikatorem
pierwszego rzędu w $M$.

Wewnątrz $M$ ustalamy
dobre uporządkowanie zbioru
$\mathbb P$ i przez rekurencję
po $n\in\omega$ wybieramy
$p_{n+1}\leq p_n$,
$p_0=p$, tak by $p_{n+1}$
ustalał jedną wartość
$\dot f(\check n)$.
Cały ciąg $(p_n)_{n<\omega}$
należy do $M$ dzięki rekurencji
i zastępowaniu w $M$.
Przeliczalna domkniętość w $M$
daje $q\leq p_n$ dla wszystkich $n$.
Każdy $p_{n+1}$ wymusza
$\dot f(\check n)=\check\gamma_n$
dla jednoznacznej $\gamma_n\in M$.
Monotoniczność zapewnia, że
$q$ wymusza wszystkie te równości.
Ciąg $(\gamma_n)_{n<\omega}$
należy do $M$ przez zastępowanie,
gdyż zbiory warunków
$p_{n+1}$ oraz wymuszane
wartości są definowalne w $M$.
Dla każdego filtru generycznego
zawierającego $q$ interpretacja
$\dot f$ jest równa temu ciągowi.

Zbiór warunków $q\leq p$,
które rozstrzygają cały ciąg
jako ciąg należący do $M$,
jest gęsty poniżej $p$ i należy
do $M$ przez definiowalność
wymuszania. Lemat o lokalnej
generyczności z Rozdziału
\ref{ch:forcing-theorem} pokazuje,
że $G$ zawiera taki warunek.
Zatem $f\in M$.

Każdy podzbiór $\omega$
jest kodowany ciągiem binarnym,
więc nie powstają nowe podzbiory
$\omega$. Gdyby $\omega_1^M$
stało się przeliczalne w $M[G]$,
istniałaby nowa funkcja
z $\omega$ na $\omega_1^M$.
Wykazaliśmy, że taka funkcja
musiałaby należeć do $M$,
co przeczy definicji $\omega_1^M$.
Ponieważ nie ma nowych liczb
porządkowych, otrzymujemy
$\omega_1^{M[G]}=\omega_1^M$.
\end{proof}

\begin{uwaga}[Różnica między domkniętością a ccc]
\label{rem:forcing34-two-properties}
Przeliczalna domkniętość i ccc
wyrażają odmienne mechanizmy.
Pierwsza pozwala uzyskać wspólne
wzmocnienie całego przeliczalnego
ciągu warunków i tym samym
rozstrzygnąć jednocześnie
przeliczalnie wiele zdań.
Druga kontroluje rozmiar zbiorów
parami niezgodnych warunków
i zapobiega zmniejszaniu liczb
kardynalnych. Nie wolno
wnioskować z samej przeliczalnej
domkniętości o zachowaniu
wszystkich liczb kardynalnych:
w następnym przykładzie
celowo zmienimy moc dużego zbioru.
\end{uwaga}

\section{Wymuszanie hipotezy continuum}

\begin{definicja}[Pojęcie wymuszania kolapsu]
\label{def:forcing34-collapse}
Niech $\lambda=(2^{\aleph_0})^M$.
Wewnątrz $M$ definiujemy
$\mathbb Q=\operatorname{Coll}(\omega_1,\lambda)$
jako zbiór wszystkich funkcji częściowych
$p:\omega_1^M\longrightarrow\lambda$,
których dziedzina jest
\emph{co najwyżej przeliczalna w $M$}.
Porządkiem jest odwrotna inkluzja,
a największym warunkiem funkcja pusta.
\end{definicja}

\begin{lemat}[Kolaps jest przeliczalnie domknięty]
\label{lem:forcing34-collapse-closed}
Pojęcie wymuszania
$\operatorname{Coll}(\omega_1,\lambda)$
z Definicji~\ref{def:forcing34-collapse}
jest przeliczalnie domknięte w $M$.
\end{lemat}
\begin{proof}
Niech $(p_n)_{n<\omega}\in M$
będzie malejącym ciągiem
warunków, czyli $p_{n+1}\supseteq p_n$.
Funkcja $q=\bigcup_{n<\omega}p_n$
jest dobrze określona, ponieważ
warunki w ciągu są zgodne.
Z aksjomatu wyboru w $M$
wynika, że suma przeliczalnie
wielu przeliczalnych dziedzin
jest przeliczalna. Stąd
$q\in\mathbb Q$, a $q\leq p_n$
dla każdego $n$.
\end{proof}

\begin{lemat}[Generyczna surjekcja kolapsu]
\label{lem:forcing34-collapse-surjection}
Jeśli $H\subseteq\mathbb Q$
jest filtrem $M$-generycznym,
to $F=\bigcup H\in M[H]$
jest surjekcją
$F:\omega_1^M\longrightarrow\lambda$.
\end{lemat}
\begin{proof}
Dla każdego $\xi<\omega_1^M$
zbiór warunków, których dziedzina
zawiera $\xi$, jest gęsty i należy
do $M$: jeśli $\xi$ nie występuje,
przedłużamy funkcję dowolną
wartością, na przykład $0$.
Dlatego $F$ jest funkcją całkowitą.

Dla każdego $\eta<\lambda$
zbiór warunków, których obraz
zawiera $\eta$, jest także gęsty.
Istotnie, dziedzina dowolnego
warunku jest przeliczalna w $M$,
a $\omega_1^M$ jest w $M$
nieprzeliczalna, więc można
wybrać nową współrzędną $\xi$
i przypisać jej $\eta$.
Generyczność zapewnia
$\eta\in\operatorname{ran}F$.
Ponieważ $\eta$ była dowolna,
$F$ jest surjekcją.
\end{proof}

\begin{twierdzenie}[Generyczny model spełniający CH]
\label{thm:forcing34-ch}
Niech $M$ będzie przeliczalnym
tranzytywnym modelem zbiorowym ZFC,
$\lambda=(2^{\aleph_0})^M$
oraz $\mathbb Q$ pojęciem wymuszania
z Definicji~\ref{def:forcing34-collapse}.
Wtedy dla każdego filtru
$M$-generycznego $H\subseteq\mathbb Q$
zachodzi
\[
 M[H]\models\mathrm{ZFC}+\mathrm{CH}.
\]
\end{twierdzenie}
\begin{proof}
Twierdzenie~\ref{thm:forcing33-zfc}
daje $M[H]\models\mathrm{ZFC}$.
Z Lematu~\ref{lem:forcing34-collapse-closed}
oraz Lematu~\ref{lem:forcing34-closed-no-sequences}
otrzymujemy
\[
 \mathcal P(\omega)^{M[H]}
 =\mathcal P(\omega)^M,
 \qquad
 \omega_1^{M[H]}=\omega_1^M.
\]
Ponieważ w $M$ istnieje bijekcja
między $\lambda$ a $\mathcal P(\omega)^M$,
taka bijekcja należy również
do $M[H]$.
Z Lematu~\ref{lem:forcing34-collapse-surjection}
istnieje tam surjekcja
$\omega_1^{M[H]}\to\lambda$.
Jej złożenie z powyższą bijekcją
jest surjekcją z $\omega_1^{M[H]}$
na wszystkie podzbiory $\omega$
w $M[H]$. Stąd
$2^{\aleph_0}\leq\aleph_1$
wewnątrz $M[H]$.
Z twierdzenia Cantora i ZFC
mamy zawsze
$2^{\aleph_0}\geq\aleph_1$.
Zatem $M[H]\models\mathrm{CH}$.
\end{proof}

\begin{przyklad}[Co zmienia kolaps]
\label{ex:forcing34-collapse-change}
Jeżeli $M\models
2^{\aleph_0}=\aleph_3$, to
$\lambda=\aleph_3^M$.
Po wymuszeniu kolapsu nadal
obowiązuje
$\mathcal P(\omega)^{M[H]}
=\mathcal P(\omega)^M$:
żaden stary ani nowy podzbiór
$\omega$ nie znika i nie
powstaje nowy. Zmienia się
jednak relacja równoliczności:
istnieje nowa surjekcja
$\omega_1^M\to\lambda$.
Liczba porządkowa $\lambda$
pozostaje taka sama, ale nie
jest już mocą $\aleph_3$
w modelu rozszerzonym.
Nie ma tu sprzeczności:
wymuszanie może zmienić
moc starego zbioru bez
zmiany jego elementów.
\end{przyklad}

\section{Wynik modelowy i granice niezależności}

\begin{twierdzenie}[Dwa modele o różnych wartościach CH]
\label{thm:forcing34-two-models}
Jeżeli istnieje przeliczalny
tranzytywny model zbiorowy
$M\models\mathrm{ZFC}$, to
istnieją przeliczalne tranzytywne
modele zbiorowe $M_0$ i $M_1$,
dla których
\[
 M_0\models\mathrm{ZFC}+\mathrm{CH},
 \qquad
 M_1\models\mathrm{ZFC}+\neg\mathrm{CH}.
\]
\end{twierdzenie}
\begin{proof}
Niech $M_0=M[H]$, gdzie
$H$ jest filtrem generycznym
dla kolapsu z
Twierdzenia~\ref{thm:forcing34-ch}.
Niech $M_1=M[G]$, gdzie
$G$ jest filtrem generycznym
dla $\operatorname{Add}(\omega,\omega_2^M)$
z Twierdzenia~\ref{thm:forcing34-not-ch}.
Istnienie filtrów zapewnia
twierdzenie Rasiowej--Sikorskiego
z Rozdziału~\ref{ch:forcing-foundations}.
Oba rozszerzenia są zbiorami
tranzytywnymi spełniającymi ZFC.
Oba rozszerzenia są także zewnętrznie
przeliczalne: każdy ich element jest
wartością pewnej nazwy $\tau\in M$,
a $M$ jest zewnętrznie przeliczalnym
zbiorem. Nie twierdzimy przy tym,
że \emph{wewnątrz} $M$ istnieje
zbiór wszystkich nazw:
nazwy tworzą tam klasę właściwą.
Korzystamy tylko z zewnętrznego
faktu, że zbiór wszystkich nazw
należących do $M$ jest podzbiorem $M$.
Twierdzenia o CH i $\neg\mathrm{CH}$
kończą dowód.
\end{proof}

\begin{uwaga}[Czego wymaga pełna niezależność]
\label{rem:forcing34-consistency}
Właśnie udowodniliśmy implikację
\emph{z istnienia przeliczalnego
tranzytywnego modelu ZFC} do istnienia
dwóch modeli o przeciwnych wartościach
CH. Nie jest to jeszcze dowód obu
implikacji metamatematycznych
\[
\begin{aligned}
 \operatorname{Con}(\mathrm{ZFC})
  &\Longrightarrow
       \operatorname{Con}(\mathrm{ZFC}+\mathrm{CH}),\\
 \operatorname{Con}(\mathrm{ZFC})
  &\Longrightarrow
       \operatorname{Con}(\mathrm{ZFC}+\neg\mathrm{CH}).
\end{aligned}
\]
Istnienie przeliczalnego
tranzytywnego modelu zbiorowego
jest mocniejszym założeniem niż
sama niesprzeczność ZFC.
Pełny wynik o względnej
niesprzeczności można uzyskać
przez formalizację argumentu
wymuszania, modele o wartościach
boolowskich lub odpowiedni
argument metamatematyczny.
Nie przedstawiamy żadnej z tych
metod jako udowodnionej bez
wykonania jej poszczególnych kroków.
Konstrukcja modeli wewnętrznych
z Rozdziału~\ref{ch:constructible}
stanowi drugi istotny kierunek
badań nad CH i GCH.
\end{uwaga}
